\begin{tikzpicture}[
    note/.style={
        draw,
        rounded corners=2pt,
        thick,
        align=center,
        inner xsep=3mm,
        inner ysep=2mm,
        font=\small
    }
]

    % Stavy automatu
    \initialstate{q0}{(0,0)}{$q_0$}{(-1.2,0)}{west}
    \state{p}{(3.0,0)}{$q$}
    \acceptingstate{r}{(6.2,0)}{$r$}

    % Přechody
    \transitionarrow[bend left=20]{q0}{above}{$0^i$}{p}
    \transitionarrow[bend right=20]{q0}{below}{$0^j$}{p}
    \transitionarrow{p}{above}{$1^i$}{r}

    % Poznámka ke stejnému stavu
    \node[note] (same) at (2.0,-2.0)
    {Po $0^i$ i $0^j$ jsme ve stejném stavu $q$};

    \draw[->,thick,dashed] (same.north) -- (p);

    % Dvě slova s různým členstvím v jazyce
    \node[note] (w1) at (4.7,-3.9) {$0^i1^i \in L_\text{eq}$};
    \node[note] (w2) at (7.7,-3.9) {$0^j1^i \notin L_\text{eq}$};

    \draw[->,thick,dashed] (r) -- (w1.north);
    \draw[->,thick,dashed] (r) -- (w2.north);

    % Spor
    \node[note,fill=red!10] (contr) at (6.2,-5.7)
    {Stejný stav $q$ a stejné pokračování $1^i$\\
     musí vést ke stejnému výsledku.\\
     To je spor.};

    % \draw[->,thick] (w1.south) -- (contr.north west);
    % \draw[->,thick] (w2.south) -- (contr.north east);

\end{tikzpicture}